\section*{Hoofdstuk \ref{ch:b2:polymer-synthesis} --- Polymeren: synthese, structuur en eigenschappen}

\begin{solution}{exo:b2:polymer-synthesis:1}
Polyetheen \ce{-[CH2-CH2]-}; polypropeen \ce{-[CH2-CH(CH3)]-}; poly(vinylchloride)
\ce{-[CH2-CHCl]-}; polystyreen \ce{-[CH2-CH(C6H5)]-}; nylon-6,6 \ce{-[NH(CH2)6NH-CO(CH2)4CO]-}.
\end{solution}

\begin{solution}{exo:b2:polymer-synthesis:2}
PET: stapgroei (polyester, water komt vrij). Polystyreen: ketengroei. Nylon-6,6: stapgroei.
Poly(methylmethacrylaat): ketengroei (een \omterm{def:b2:polymer-synthesis:polymer}{monomeer} met \ce{C=C}). Poly(melkzuur) uit het hydroxyzuur:
stapgroei, een \omterm{def:b2:polymer-synthesis:polymer}{monomeer} \ce{A-B}.
\end{solution}

\begin{solution}{exo:b2:polymer-synthesis:3}
$\bar M_n = 1000 \times \qty{104.15}{g/mol} \approx \qty{104}{kg/mol}$ (styreen \ce{C8H8},
\qty{104.152}{g/mol}).
% ledger: aw:C, aw:H
\end{solution}

\begin{solution}{exo:b2:polymer-synthesis:4}
\omterm{def:b2:polymer-synthesis:tacticity}{Syndiotactisch}. Ja: een regelmatige keten kan zich tot kristallieten schikken.
\end{solution}

\begin{solution}{exo:b2:polymer-synthesis:5}
Massafracties 0.5 en 0.5. $1/\bar M_n = 0.5/10 + 0.5/100 = 0.055$, $\bar M_n \approx
\qty{18.2}{kg/mol}$; $\bar M_w = 0.5 \times 10 + 0.5 \times 100 = \qty{55}{kg/mol}$; $\text{\DH} \approx
3.0$.
\end{solution}

\begin{solution}{exo:b2:polymer-synthesis:6}
$1/(1-0.98) = 50$; $1/(1-0.995) = 200$.
\end{solution}

\begin{solution}{exo:b2:polymer-synthesis:7}
$r = 1/1.01$; bij $p = 1$ is $\bar X_n = (1+r)/(1-r) = (1.01 + 1)/(1.01 - 1) = 201$.
\end{solution}

\begin{solution}{exo:b2:polymer-synthesis:8}
$R_p \propto [\mathrm I]^{1/2}$: vermenigvuldigd met $\sqrt2 \approx 1.41$. $\nu \propto [\mathrm I]^{-1/2}$:
gedeeld door $\sqrt 2$, ongeveer 0.71 keer zijn vroegere waarde.
\end{solution}

\begin{solution}{exo:b2:polymer-synthesis:9}
Zuiver pvc heeft zijn glasovergang boven kamertemperatuur: een star glas. De kleine
weekmakermoleculen zitten tussen de ketens, houden ze uit elkaar en laten segmenten bij lagere temperatuur bewegen: $T_g$ daalt
onder kamertemperatuur, en het materiaal is bij gebruik rubberachtig en buigzaam.
\end{solution}

\begin{solution}{exo:b2:polymer-synthesis:10}
$\dfrac{\mathrm d}{\mathrm dx}\bigl(x p^{\,x-1}\bigr) = p^{\,x-1}(1 + x\ln p)$, nul voor $x = -1/\ln p$
(een maximum: het teken gaat van $+$ naar $-$). Voor $p$ dicht bij 1 is $\ln p \approx -(1-p)$, dus $x \approx
1/(1-p) = \bar X_n$.
\end{solution}

\begin{solution}{exo:b2:polymer-synthesis:11}
Fase 1: \ce{C6H4(COOCH3)2 + 2HOCH2CH2OH -> C6H4(COOCH2CH2OH)2 + 2CH3OH}, een \omterm{def:b2:acyl-substitution:transesterification}{omestering}
waarbij het methanol afgedestilleerd wordt. Fase 2: elk diestermolecuul draagt twee uiteinden \ce{CH2CH2OH}; een uiteinde valt een ester
van een ander molecuul aan en maakt een molecuul ethaan-1,2-diol vrij, dat weggepompt wordt. Het \omterm{def:b2:polymer-synthesis:polymer}{monomeer}
draagt beide soorten groepen in de juiste verhouding, zoals een \omterm{def:b2:polymer-synthesis:polymer}{monomeer} \ce{A-B}: de balans $r = 1$ zit
ingebouwd, en het verwijderen van het diol drijft $p$ naar 1.
\end{solution}

\begin{solution}{exo:b2:polymer-synthesis:12}
Equimolaire voeding: verhouding $= (2.0 + 1)/(1 + 0.5) = 2$: twee keer zoveel eenheden van \omterm{def:b2:polymer-synthesis:polymer}{monomeer} 1. Met $r_1 = r_2 =
0$ is de verhouding $= [\mathrm M_1][\mathrm M_2]/([\mathrm M_2][\mathrm M_1]) = 1$, ongeacht de voeding: elk
radicaal addeert alleen het andere \omterm{def:b2:polymer-synthesis:polymer}{monomeer}, een alternerend \omterm{def:b2:polymer-synthesis:copolymer}{copolymeer}.
\end{solution}

\begin{solution}{pb:b2:polymer-synthesis:1}
\textbf{1.} \ce{H2N(CH2)6NH2} en \ce{HOOC(CH2)4COOH}.
\textbf{2.} Eén water per binding:
\[\mathrm{{-}COOH + H_2N{-} \longrightarrow {-}CO{-}NH{-} + H_2O}.\]
\textbf{3.} Het zout bevat precies één diamine per dizuur; afzonderlijk afwegen
bereikt nooit de exactheid die de vergelijking van Carothers vraagt.
\textbf{4.} \ce{-[NH(CH2)6NH-CO(CH2)4CO]-}, \ce{C12H22N2O2}, \qty{226.32}{g/mol}.
\textbf{5.} De \omterm{def:b2:polymer-synthesis:polymer}{repeterende eenheid} bevat twee eenheden die van een \omterm{def:b2:polymer-synthesis:polymer}{monomeer} afkomstig zijn: $M_0 = 226.32/2 = \qty{113.16}{g/mol}$.
\textbf{6.} De koolstofatomen van het diamine (6) en van het dizuur (6).
\textbf{7.} $\bar X_n = 50$, $\bar M_n = 50 \times 113.16 \approx \qty{5.66}{kg/mol}$.
\textbf{8.} $\bar X_n = 100$, $\bar M_n \approx \qty{11.3}{kg/mol}$.
\textbf{9.} Als 99~\% van de groepen gereageerd heeft, heeft de gemiddelde keten maar 100 eenheden, te kort voor een
sterke vezel; elke verdere stap naar volledige omzetting verdubbelt of verdrievoudigt de ketenlengte.
\textbf{10.} $\bar X_n = 10\,000/113.16 = 88.4$; $p = 1 - 1/88.4 \approx 0.989$.
\textbf{11.} Door water te verwijderen: de smelt wordt ver boven haar smeltpunt verhit, op het einde onder verlaagde
druk, zodat het evenwicht naar het amide blijft verschuiven.
\textbf{12.} Een monofunctionele onzuiverheid sluit ketenuiteinden af, en elke onzuiverheid verstoort de balans 1\,:\,1.
\textbf{13.} $r = 1/1.01 \approx 0.990$.
\textbf{14.} $\bar X_n = (1+r)/(1-r) = 201$; $\bar M_n \approx 201 \times 113.16 \approx
\qty{22.7}{kg/mol}$.
\textbf{15.} $\bar X_n = 1.9901/(1.9901 - 2 \times 0.9901 \times 0.99) \approx 1.9901/0.0297 \approx
67$.
\textbf{16.} Aminegroepen: wanneer de zuurgroepen opgebruikt zijn, eindigt elke keten op \ce{NH2}.
\textbf{17.} Het zet een amine-uiteinde om in een amide dat niet verder kan reageren: het sluit ketens af, beperkt zo
$\bar M_n$ en stopt verdere groei wanneer het \omterm{def:b2:polymer-synthesis:polymer}{polymeer} opnieuw gesmolten wordt.
\textbf{18.} Om een reproduceerbare smeltviscositeit te behouden voor het spinnen en vormen, die anders
zou verlopen doordat ketens in de smelt blijven reageren.
\textbf{19.} $\text{\DH} = 1 + p = 1.99$.
\textbf{20.} $\bar M_w = 1.99 \times 11.3 \approx \qty{22.5}{kg/mol}$.
\textbf{21.} Aantalfractie $n_1 = 1 - p = 0.01$, 1~\% van de moleculen; massafractie $w_1 = (1-p)^2 =
10^{-4}$, 0.01~\% van de massa.
\textbf{22.} Rond $x = -1/\ln 0.99 \approx 99.5$, dus rond $\bar X_n = 100$ (oefening 10).
\textbf{23.} $1000/226.32 = \qty{4.42}{mol}$ \omterm{def:b2:polymer-synthesis:polymer}{repeterende eenheden}, met elk twee amidebindingen: \qty{8.84}{mol}
water, $8.84 \times 18.015 \approx \qty{159}{g}$.
\textbf{24.} Te korte ketens geven een zwakke vezel; te lange ketens geven een smelt die te viskeus is om
door de fijne gaatjes van de spindop geperst te worden.
\textbf{25.} $\bar X_n = 20\,000/113.16 = 176.7$ en $p = 1 - 1/176.7 = \boldsymbol{\approx 0.994}$.
\end{solution}
